\begin{statement}{A.2.4}{Proved}{Pinch points and cuspidal edges}
At the pinch point $F=x^2-y^2z$, with normalization $x=uv,y=u,z=v^2$, write $B=\C\{u,v\}$. Then
\[
 K=B\,dF,\qquad T\simeq B/R\simeq R/(x,y),
\quad \tau=2yz\,dx-2xz\,dy-xy\,dz
\tag{A.8}
\]
is a torsion generator. At the cuspidal edge $F=x^2-y^3$, with normalization $x=t^3,y=t^2$ and unchanged coordinate $z$, write $B=\C\{t,z\}$. Then
\[
 K=t^{-1}B\,dF,\qquad
 T\simeq t^{-1}B/R\simeq R/(x,y^2),
 \quad\tau=2y\,dx-3x\,dy.
\tag{A.9}
\]
The conductor is $(x,y)$ for the pinch point and $(x,y)$ for the cusp; in the cusp normalization it is $t^2B$. An ordinary surface, meaning a surface whose singularities are double normal crossings, triple normal crossings, and pinch points, satisfies canonically $K=\eta_*\eta^*N^*$ and $T=N^*\otimes(\eta_*\OO_{\widetilde S}/\OO_S)$. For the cuspidal edge the additional quotient $J^{-1}/B$ is generated by $t^{-1}$ over $\C\{z\}$.
\end{statement}
\begin{proof}
For the pinch point,
\[
 R=\C\{u,v^2,uv\}=\C\{v^2\}+uB,
 \quad J=(uv,uv^2,u^2)=u(u,v)B.
\]
The equality of the two descriptions of $J$ follows by separating even and odd powers of $v$. An element $q\in J^{-1}$ satisfies $(qu)(u,v)\subset B$. Since $B$ is a regular local ring, $\Hom_B((u,v),B)=B$: a meromorphic function whose products with $u$ and $v$ are holomorphic has no pole along any divisor, and its possible isolated singularity is removable by the holomorphic extension theorem. Hence $qu=h\in B$. The requirements $hv,hv^2\in R$, restricted to $u=0$, say respectively that $h(0,v)v$ and $h(0,v)v^2$ are even power series. The first requires $h(0,v)$ to be odd and the second to be even. Therefore $h(0,v)=0$, so $h\in uB$ and $q\in B$. The reverse inclusion follows either directly or from (A.3). Also $B/R$ is generated by $v$ and its annihilator is $(u,uv)=(y,x)$. Multiplying $dF$ by $v=x/y$ gives the displayed $\tau$.

For the cusp, $R=\C\{t^2,t^3,z\}$ and $J=(t^3,t^4)$. A meromorphic $q$ with $qt^3,qt^4\in R$ has no poles away from $t=0$; writing its Laurent expansion in $t$ shows that the lowest allowed exponent is $-1$. Indeed an exponent $-3$ produces a forbidden $t$ term after multiplication by $t^4$, an exponent $-2$ produces one after multiplication by $t^3$, and smaller exponents produce a pole. Conversely every $t^{-1}B$ has both products in $R$. Thus $J^{-1}=t^{-1}B$. Modulo $R$, the quotient has the $\C\{z\}$-basis $t^{-1},t$; multiplication by $y$ sends the first basis vector to the second, while $x$ and $y^2$ kill both. This proves the cyclic presentation and (A.9). The conductor calculations follow directly from the displayed descriptions of $R\subset B$. Formula (A.2) and the three local computations prove the ordinary-surface assertion by gluing.
\end{proof}


% A.2.5 from part_A.tex

